% fig14_scaling_laws.tex — 图14 Scaling Law 示意曲线
% 《深度学习与大语言模型：前沿技术全景报告》图示库
% 由 dl_llm_report.tex 抽取，经 % fig14_scaling_laws.tex — 图14 Scaling Law 示意曲线
% 《深度学习与大语言模型：前沿技术全景报告》图示库
% 由 dl_llm_report.tex 抽取，经 % fig14_scaling_laws.tex — 图14 Scaling Law 示意曲线
% 《深度学习与大语言模型：前沿技术全景报告》图示库
% 由 dl_llm_report.tex 抽取，经 % fig14_scaling_laws.tex — 图14 Scaling Law 示意曲线
% 《深度学习与大语言模型：前沿技术全景报告》图示库
% 由 dl_llm_report.tex 抽取，经 \input{figs/fig14_scaling_laws} 引用（caption/label 在主文件）
% 注意：本文件为片段，依赖主文件导言区的 tikz/pgfplots 宏包、颜色定义与 tikzset 样式，不可单独编译。

\begin{tikzpicture}
\begin{semilogxaxis}[width=0.62\textwidth, height=6.4cm,
  xlabel={训练算力（对数刻度）}, ylabel={验证损失（示意）},
  xtick=\empty, ytick=\empty, axis lines=left, ymin=0.8, ymax=3.8,
  legend pos=north east, legend style={font=\scriptsize, draw=cgray, fill=white}]
\addplot[very thick, cgray] coordinates {(1,3.4)(3,2.8)(10,2.35)(30,2.0)(100,1.72)(300,1.5)(1000,1.32)};
\addlegendentry{参数优先（Kaplan 路线）}
\addplot[very thick, cblue] coordinates {(1,2.95)(3,2.4)(10,2.0)(30,1.7)(100,1.46)(300,1.27)(1000,1.12)};
\addlegendentry{数据等比（Chinchilla）}
\addplot[thick, cgreen!70!black, dashed] coordinates {(1,2.7)(3,2.15)(10,1.78)(30,1.5)(100,1.28)(300,1.11)(1000,0.98)};
\addlegendentry{+ 高质量/合成数据}
\node[lab, anchor=south west] at (axis cs:30,2.42) {同算力下损失更低};
\node[lab, anchor=west] at (axis cs:10,0.95) {幂律：损失随算力对数线性下降};
\end{semilogxaxis}
\end{tikzpicture}
 引用（caption/label 在主文件）
% 注意：本文件为片段，依赖主文件导言区的 tikz/pgfplots 宏包、颜色定义与 tikzset 样式，不可单独编译。

\begin{tikzpicture}
\begin{semilogxaxis}[width=0.62\textwidth, height=6.4cm,
  xlabel={训练算力（对数刻度）}, ylabel={验证损失（示意）},
  xtick=\empty, ytick=\empty, axis lines=left, ymin=0.8, ymax=3.8,
  legend pos=north east, legend style={font=\scriptsize, draw=cgray, fill=white}]
\addplot[very thick, cgray] coordinates {(1,3.4)(3,2.8)(10,2.35)(30,2.0)(100,1.72)(300,1.5)(1000,1.32)};
\addlegendentry{参数优先（Kaplan 路线）}
\addplot[very thick, cblue] coordinates {(1,2.95)(3,2.4)(10,2.0)(30,1.7)(100,1.46)(300,1.27)(1000,1.12)};
\addlegendentry{数据等比（Chinchilla）}
\addplot[thick, cgreen!70!black, dashed] coordinates {(1,2.7)(3,2.15)(10,1.78)(30,1.5)(100,1.28)(300,1.11)(1000,0.98)};
\addlegendentry{+ 高质量/合成数据}
\node[lab, anchor=south west] at (axis cs:30,2.42) {同算力下损失更低};
\node[lab, anchor=west] at (axis cs:10,0.95) {幂律：损失随算力对数线性下降};
\end{semilogxaxis}
\end{tikzpicture}
 引用（caption/label 在主文件）
% 注意：本文件为片段，依赖主文件导言区的 tikz/pgfplots 宏包、颜色定义与 tikzset 样式，不可单独编译。

\begin{tikzpicture}
\begin{semilogxaxis}[width=0.62\textwidth, height=6.4cm,
  xlabel={训练算力（对数刻度）}, ylabel={验证损失（示意）},
  xtick=\empty, ytick=\empty, axis lines=left, ymin=0.8, ymax=3.8,
  legend pos=north east, legend style={font=\scriptsize, draw=cgray, fill=white}]
\addplot[very thick, cgray] coordinates {(1,3.4)(3,2.8)(10,2.35)(30,2.0)(100,1.72)(300,1.5)(1000,1.32)};
\addlegendentry{参数优先（Kaplan 路线）}
\addplot[very thick, cblue] coordinates {(1,2.95)(3,2.4)(10,2.0)(30,1.7)(100,1.46)(300,1.27)(1000,1.12)};
\addlegendentry{数据等比（Chinchilla）}
\addplot[thick, cgreen!70!black, dashed] coordinates {(1,2.7)(3,2.15)(10,1.78)(30,1.5)(100,1.28)(300,1.11)(1000,0.98)};
\addlegendentry{+ 高质量/合成数据}
\node[lab, anchor=south west] at (axis cs:30,2.42) {同算力下损失更低};
\node[lab, anchor=west] at (axis cs:10,0.95) {幂律：损失随算力对数线性下降};
\end{semilogxaxis}
\end{tikzpicture}
 引用（caption/label 在主文件）
% 注意：本文件为片段，依赖主文件导言区的 tikz/pgfplots 宏包、颜色定义与 tikzset 样式，不可单独编译。

\begin{tikzpicture}
\begin{semilogxaxis}[width=0.62\textwidth, height=6.4cm,
  xlabel={训练算力（对数刻度）}, ylabel={验证损失（示意）},
  xtick=\empty, ytick=\empty, axis lines=left, ymin=0.8, ymax=3.8,
  legend pos=north east, legend style={font=\scriptsize, draw=cgray, fill=white}]
\addplot[very thick, cgray] coordinates {(1,3.4)(3,2.8)(10,2.35)(30,2.0)(100,1.72)(300,1.5)(1000,1.32)};
\addlegendentry{参数优先（Kaplan 路线）}
\addplot[very thick, cblue] coordinates {(1,2.95)(3,2.4)(10,2.0)(30,1.7)(100,1.46)(300,1.27)(1000,1.12)};
\addlegendentry{数据等比（Chinchilla）}
\addplot[thick, cgreen!70!black, dashed] coordinates {(1,2.7)(3,2.15)(10,1.78)(30,1.5)(100,1.28)(300,1.11)(1000,0.98)};
\addlegendentry{+ 高质量/合成数据}
\node[lab, anchor=south west] at (axis cs:30,2.42) {同算力下损失更低};
\node[lab, anchor=west] at (axis cs:10,0.95) {幂律：损失随算力对数线性下降};
\end{semilogxaxis}
\end{tikzpicture}
